\section{Experimental Setup}
\label{sec:setup}

We evaluate BanditGPT on a canonical portfolio of three LLMs, spanning a large cost differential, and using realistic prompt data to simulate both stationary and non-stationary production environments.

\subsection{Datasets and Ground Truth}

We use a dataset of 1,785 prompts sampled from the LMSYS Chatbot Arena~\cite{zheng2023judging}. This dataset reflects real-world, open-ended user interactions spanning diverse domains, including coding, reasoning, and conversational tasks. 

To ensure robust and reproducible evaluation, we employ an automated LLM-as-a-judge paradigm using DeepSeek-R1~\cite{deepseekai2025deepseekr1} as the evaluator.  Each response is scored on three continuous dimensions---\emph{Reasoning Quality} (40\%), \emph{Instruction Following} (30\%), and \emph{Communication Quality} (30\%)---yielding a composite scalar reward $r_t \in [0, 1]$.  This continuous rubric provides the granular quality distinctions necessary for multi-arm routing; binary win/loss outcomes collapse into ties on a large fraction of prompts and are insufficient for $K > 2$ portfolios.

For online learning, the bandit's convergence rate scales directly with the reward gap between the optimal and suboptimal arms~\cite{lattimore2020bandit}; we therefore select the judge that maximises routing margins.  DeepSeek-R1 produces the largest inter-model gaps and the highest fraction of actionable prompts among the judges we evaluated, while its oracle routing decisions capture $\geq 97\%$ of two independent supplementary judges' oracle reward (Appendix~\ref{sec:judge_robustness}).  Recent multi-judge aggregation methods~\cite{verga2024poll,zhao2025care} target confounder-induced bias rather than the signal compression inherent in averaging judges with different scale ranges; applied to our panel, averaging reduces routing margins by ${\sim}30\%$ (Appendix, Figure~\ref{fig:panel_signal}), making the single maximally discriminative judge the principled choice for bandit learning.

The 1,785 prompts are split into three sets:
1. \textbf{Warmup Set ($n=1,028$)}: Used offline to train the `Warmup Priors` that initialize the bandit's sufficient statistics ($A_a, b_a$), addressing the cold-start problem.
2. \textbf{Validation Set ($n=1,785$)}: Used for online learning and hyperparameter tuning in stationary conditions.
3. \textbf{Test Sets ($n \approx 1,800$)}: Constructed to simulate specific non-stationary scenarios (e.g., reward drift and cost drift), evaluated across multiple random seeds.

\subsection{Environment Construction}

Our portfolio ($K=3$) spans three distinct cost tiers: Llama-3.1-8B (Budget, \$0.00003/req), Mistral-Large-2512 (Mid-tier, \$0.003/req), and Gemini-2.5-Pro (Frontier, \$0.015/req). This $500\times$ price differential stresses the router's ability to balance quality and cost efficiently.

We construct two distinct non-stationary test environments to evaluate the router's robustness:
1. \textbf{Reward Drift (Model Quality Shift)}: Simulates a silent API update where model quality changes mid-deployment. The test trajectory ($n=1,785$) is split into two phases. Phase 1 (steps 1--893) uses normal rewards where Mistral-Large is the utility-optimal model. In Phase 2 (steps 894--1,785), Llama-8B and Mistral-Large exchange their reward distributions, simulating a scenario where Mistral degrades to the worst option while Llama becomes the best. Gemini-Pro serves as an unchanged anchor.
2. \textbf{Cost Drift (Price Drop)}: Simulates an API price cut. The test trajectory ($n=1,824$) is similarly split. Phase 1 (steps 1--912) uses normal pricing. In Phase 2 (steps 913--1,824), Gemini-2.5-Pro's price drops to \$0.0001 per request, drastically altering the optimal cost-quality trade-off.

\subsection{Baselines}

We compare BanditGPT against two standard baselines representing different levels of routing sophistication:
1. \textbf{Fixed Policy (Offline):} The router is initialized with warmup priors but performs zero online learning ($\gamma=1.0$, updates disabled). A static, offline-tuned cost penalty $\lambda_s$ is applied. This represents the industry standard of deploying a frozen classifier trained on offline data.
2. \textbf{Naive Bandit ($\gamma=1.0$):} A standard LinUCB algorithm initialized with warmup priors. It learns online but retains infinite memory (no forgetting). A static cost penalty $\lambda_s$ is applied. This represents a straightforward application of online learning without adaptation for non-stationarity or budget pacing.
3. \textbf{BanditGPT ($\gamma=0.997$):} Our proposed system, initialized with warmup priors, utilizing geometric forgetting, and governed by the Primal-Dual `BudgetPacer` targeting a specific dollar budget.

All three conditions use the same exploration rate ($\alpha=0.1$) and prior strength ($n_{eff}=10$).

\subsection{Metrics}

We evaluate the systems using three primary metrics:
1. \textbf{Cumulative Regret:} Measured against a cost-adjusted oracle that selects the model maximizing $r_{a,t} - \lambda_s \tilde{c}_a$ at each step $t$ (where $\lambda_s$ is a shared static penalty). Lower regret indicates better routing decisions relative to the optimal choice.
2. \textbf{Budget Compliance:} The ratio of actual realized cost to the target budget $B$. A perfect pacer achieves a ratio $\leq 1.0\times$, maximizing quality without overspending. Ratios $>1.0\times$ indicate budget violations.
3. \textbf{Pareto AUC:} The area under the quality-cost Pareto frontier, measuring the router's efficiency across all possible cost trade-offs. Higher AUC indicates better quality at lower costs.